% Chapter 1, typeset from lectures/L01: the README is the opener and the closing summary, and
% each appendix is a section, in the same order and with the same subsections.

\chapter{The AVR Core and the Toolchain}
\label{c1:ch:core}

{\large\itshape What the machine is, and how to make it run something.\par}

\begin{chapterbox}{In this chapter}
\begin{itemize}
  \item What kind of machine an ATmega328P is, and which of your C intuitions survive contact
    with it.
  \item The 32 registers, and why only half of them can be loaded with a constant.
  \item \code{SREG}: the flags every arithmetic instruction writes and every branch reads.
  \item Three separate address spaces, and why address \code{0x0060} names three different bytes.
  \item What an instruction \emph{is}: \code{ldi r16, 0x2A} taken apart into the sixteen bits it
    assembles to.
  \item The toolchain: \code{avra} as an assembler, its listing and map files, and a simulator
    that runs your code and counts its cycles.
  \item Your first subroutine, \code{shift_bits}, costed on paper before it is run.
  \item A program that lights an LED in two instructions, which is as far as one pin gets you.
\end{itemize}
\end{chapterbox}

The chapter is worked in the order a first session at the machine should go. The machine comes
first, before any code: registers, flags and the three address spaces, read off the figures in
\secref{c1:sec:core}. Nothing runs yet; the point is to know what you are writing \emph{for}. Then
one instruction, taken apart, and the toolchain end to end: one \code{.asm} file through
\code{avra} to a hex, then the same hex into a disassembler and into a simulator, three views of
one artifact. Then counting cycles, and finally writing something, so that the chapter ends with
something you can see.

Two predictions are worth making before anything runs:
\begin{itemize}
  \item what \code{ldi r16, 0x2A} assembles to, from the encoding alone.
  \item how many cycles \code{shift_bits} takes for an argument of 5, from the instruction table
    alone.
\end{itemize}
You will check both, and one of them is the subject of this chapter's cross-check exercise.

% Section 1.1, from lectures/L01/appendix/a_avr_core.md.

\section{The AVR Core}
\label{c1:sec:core}

\subsection{What kind of machine this is}
\label{c1:sec:machine}
An ATmega328P is an 8-bit microcontroller with 32 registers, 2\,KB of RAM, 32\,KB of flash, and
no operating system, no memory protection, and no cache. It runs one instruction at a time, and
almost every instruction takes one clock cycle.

Three of its properties will contradict something you know from C, and it is worth naming them
before they surprise you rather than after.

\textbf{It is a load-store machine.} Arithmetic happens between registers and nowhere else. There
is no instruction that adds two bytes of memory together; you load one into a register, load the
other into another register, add, and store the result back. In C you write \code{a += b} and the
compiler decides how many instructions that is. Here you are the compiler.

\textbf{It is a Harvard machine.} Program memory and data memory are separate address spaces with
separate instructions. A pointer to a constant string in flash and a pointer to a variable in RAM
are both 16-bit numbers, and they are not interchangeable, and nothing checks. This is the single
biggest departure from the flat address space C leads you to expect (\secref{c1:sec:spaces}).

\textbf{It has 32 registers, and they are not all equal.} They are equal to most instructions and
unequal to a handful of important ones, and the handful is what shapes real AVR code
(\secref{c1:sec:regfile}).

What it is \emph{not} is slow or crude. At 16\,MHz it executes 16 million instructions a second,
and almost every one of them does something useful. The reason to write assembly here is not that
the compiler is bad; it is that on a machine this small you can hold the whole of it in your head,
and knowing exactly what a routine costs is worth something.

\subsection{The register file}
\label{c1:sec:regfile}
Thirty-two 8-bit registers, named \code{r0} to \code{r31}. Every one of them can be a source or a
destination for \code{mov}, \code{add}, \code{and}, \code{or}, \code{ld}, \code{st} and most of the
rest of the instruction set.

\bookfigure{register_file}{The 32 AVR registers as two columns of sixteen. The upper half is
the half immediate instructions can reach, and \code{r26} to \code{r31} are the pointer pairs X,
Y and Z.}{fig:regfile}

Two things about that picture matter more than the other thirty.

\textbf{Only \code{r16} to \code{r31} can be loaded with a constant.} \code{ldi r16, 0x2A} is
legal; \code{ldi r15, 0x2A} is not, and the assembler will tell you so. The same restriction
applies to \code{andi}, \code{ori}, \code{subi}, \code{cpi} and \code{sbci}. The reason is not
arbitrary: those instructions spend eight of their sixteen bits on the constant, which leaves
four bits for the register number, and four bits can only count to sixteen
(\secref{c1:sec:encoding}).

The practical consequence is that the upper half is where the work happens and the lower half is
where you park things. A routine that needs five working registers and reaches for \code{r10}
will find it cannot put a constant in it, and will end up loading the constant into \code{r16} and
moving it, which costs an extra instruction every time.

\textbf{\code{r26} to \code{r31} are also the three pointer registers.} Taken in pairs they form
the 16-bit pointers \code{X} (\code{r27:r26}), \code{Y} (\code{r29:r28}) and \code{Z}
(\code{r31:r30}). They are still ordinary registers; nothing stops you using \code{r30} to hold a
loop counter. But the moment you want to read memory through a pointer, they are the only six
that can do it, and a routine that has casually filled them with something else has to spill them
first. Chapter~4 is about this in earnest.

\code{r0} and \code{r1} deserve a note. \code{r0} is where \code{lpm} deposits the byte it read
from program memory, so a routine using \code{lpm} cannot also be keeping something in \code{r0}.
\code{r1} has no hardware role at all, but the C compiler keeps zero in it permanently, which
matters the moment you call C from assembly or the reverse. That is Chapter~6's problem; for now,
leaving both alone costs you nothing.

\subsection{The status register}
\label{c1:sec:sreg}
\code{SREG} is one 8-bit register holding the flags that arithmetic writes and branches read.

\bookfigure{sreg}{\code{SREG} as its eight bits. Bit 7, \code{I}, is the global interrupt
enable; bits 1 and 0, \code{Z} and \code{C}, are the ones the conditional branches
read.}{fig:sreg}

\begin{narrowtable}{@{}clp{19em}@{}}
  \thead{Bit} & \thead{Name} & \thead{Set when} \\
  \midrule
  7 & \code{I} & Interrupts are globally enabled. Set by \code{sei}, cleared by \code{cli}. \\
  6 & \code{T} & A bit copied by \code{bst}, to be copied back by \code{bld}. Rarely used. \\
  5 & \code{H} & A carry out of bit 3, for BCD arithmetic. Rarely used. \\
  4 & \code{S} & \code{N} exclusive-or \code{V}: the sign of a \emph{signed} comparison. \\
  3 & \code{V} & A signed overflow occurred. \\
  2 & \code{N} & The result's most significant bit was set. \\
  1 & \code{Z} & The result was zero. \\
  0 & \code{C} & A carry out of bit 7, or a borrow. \\
\end{narrowtable}

In practice you will use two of them constantly and the rest almost never. \code{Z} and \code{C}
are what \code{cp}, \code{cpi} and every arithmetic instruction leave behind, and what the
branches test:

\begin{narrowtable}{@{}lll@{}}
  \thead{Branch} & \thead{Taken when} & \thead{Reads} \\
  \midrule
  \code{breq} & the two values were equal & \code{Z} set \\
  \code{brne} & they were not equal & \code{Z} clear \\
  \code{brlo} & the first was lower, unsigned & \code{C} set \\
  \code{brsh} & the first was the same or higher, unsigned & \code{C} clear \\
\end{narrowtable}

\textbf{\code{cp} and \code{cpi} are subtractions whose result is thrown away.} That is the whole
of what a comparison is on this machine: \code{cp r19, r24} computes \code{r19 - r24}, discards
the difference, and keeps the flags. So \code{cp} followed by \code{breq} means ``branch if
equal'', and \code{cp} followed by \code{brlo} means ``branch if the first was lower''. If you
find yourself unsure which way round a comparison reads, write out the subtraction and the
ambiguity disappears.

\textbf{The flags are the most volatile state on the machine.} Nearly every instruction writes
some of them, so \code{cp} and its branch must be adjacent, or nearly so. A \code{mov} between
them is harmless because \code{mov} writes no flags; an \code{inc} between them is a bug, and one
that shows up as a branch going the wrong way for no visible reason. When you cannot remember
whether an instruction touches the flags, the instruction set summary has a column for it, and
that column is the reason it exists.

Bit 7, \code{I}, is different in kind from the other seven. It is not a result of anything; it is
a switch, and it controls whether interrupts can happen at all. Chapter~3 is where it matters.

\subsection{Three address spaces}
\label{c1:sec:spaces}
This is the part that most often goes wrong for a reader arriving from C, so it is worth being
slow about.

\bookfigure{memory_spaces}{Three separate address maps side by side: program memory in word
addresses, data memory from the register file up to internal SRAM, and EEPROM on its
own.}{fig:spaces}

There is no single address space on this machine. There are three, they overlap numerically, and
nothing in the instruction set or the assembler will stop you reading one when you meant another.

\textbf{Program memory} holds your instructions and any constants you put there deliberately. It
is addressed in \textbf{words}, not bytes, which is why \code{PCINT0}'s slot is quoted as
\code{0x06} while GNU \code{as} wants \code{0x0C}. The entries being 2 apart is a separate fact:
each slot is two words, so that a two-word \code{jmp} fits in one (\S3.1.2).

Only \code{lpm} reads it. An ordinary \code{ld} pointed at ``address 0x0100'' reads SRAM, not the
code at word \code{0x0100}, and does so silently.

\textbf{Data memory} is one flat space containing four things in order: the register file, the 64
I/O registers, 160 extended I/O registers, and finally 2048 bytes of actual SRAM starting at
\code{0x0100}. \code{RAMEND} is \code{0x08FF}, and it is where you point the stack
(\secref{c1:sec:program}).

The register file really is addressable there. \code{lds r16, 0x0005} loads the contents of
\code{r5}. This is true, occasionally useful for a debugger, and something no ordinary program
does, because \code{mov r16, r5} does the same thing in one cycle instead of two.

\textbf{EEPROM} is 1024 bytes of memory that survives power being removed. It has its own address
space and \textbf{no load or store instruction reaches it at all}; you write an address into
\code{EEAR}, a byte into \code{EEDR}, and set a bit in \code{EECR}. This book does not use it, but
knowing it is a third space rather than a region of the second is the point.

\subsection{The I/O window: two names for one register}
\label{c1:sec:iowindow}
Every I/O register has two addresses, and both appear in the datasheet.

\code{PORTB} is \textbf{I/O address \code{0x05}} and \textbf{data space address \code{0x25}}.
They differ by \code{0x20}, because the I/O registers begin \code{0x20} bytes into the data
space. Which number you use depends entirely on which instruction you are writing:

\begin{booktable}{@{}p{5.6em}p{9.6em}Lp{3.4em}@{}}
  \thead{Instruction} & \thead{Takes} & \thead{Reaches} & \thead{Cost} \\
  \midrule
  \code{in} / \code{out} & an I/O address, 0 to 63 & the low 64 I/O registers & 1 cycle \\
  \code{sbi} / \code{cbi} & an I/O address, 0 to 31 & the low 32, one bit at a time & 2 cycles \\
  \code{lds} / \code{sts} & a data space address & anywhere in data memory & 2 cycles \\
  \code{ld} / \code{st} & a pointer register & anywhere in data memory & 2 cycles \\
\end{booktable}

Giving one of them the other's number is the classic first AVR bug. \code{out 0x25, r16} writes
to I/O register \code{0x25}, which is data space \code{0x45}, which is \code{TCCR0B}: Timer0's
clock select and its waveform mode. Your LED does nothing, your program looks correct, and
somewhere else a timer you have not written yet is running at a rate you did not choose.

In this book you will rarely write either number, because \code{.include "m328Pdef.inc"} defines
\code{PORTB} for you, and it defines it as the \textbf{I/O} address, \code{0x05}. So
\code{out PORTB, r16} is right, and reaching the same register with \code{sts} means writing the
\code{0x20} yourself:

\begin{asm}
out PORTB, r16              ; the I/O address, 0x05
sts PORTB + 0x20, r16       ; the same register through the data space, 0x25
\end{asm}

\textbf{Not every register works that way, and the device file tells you which.} \code{WDTCSR} is
\code{0x60} and \code{TCCR1B} is \code{0x81}, data space addresses already, because those
registers live above the I/O window and \code{out} cannot reach them at all.
\code{m328Pdef.inc} marks them \code{; MEMORY MAPPED}, and that comment is the whole rule: if it
is there, the number is a data space address and you add nothing; if it is not, the number is an
I/O address and \code{sts} needs \code{+ 0x20}.

\textbf{The mistake this invites is worse than a wrong register.} Write \code{sts PORTB, r16} and
it assembles without complaint to data space \code{0x0005}, which is not a peripheral at all, but
the general-purpose register \textbf{r5}. Your port never changes, a register you were using
quietly takes the value instead, and nothing anywhere reports a problem. Assembly written for
\code{avr-gcc} does exactly this, because there \code{PORTB} already means \code{0x25}; carrying
such a line across unchanged is the single most productive source of this bug.
\secref{c1:sec:studio} has the rest of that comparison.

\subsection{What an instruction is}
\label{c1:sec:encoding}
Every AVR instruction is one or two 16-bit words. \code{ldi} is one word, and every bit of it is
accounted for.

\bookfigure{ldi_encoding}{\code{ldi r16, 0x2A} as sixteen bit cells: the pattern
\code{1110 KKKK dddd KKKK} above the assembled word, with the opcode, the two halves of the
immediate and the register field braced separately.}{fig:ldi}

\[
  \texttt{ldi }R_d, K \;\longrightarrow\; \texttt{1110}\;K_{7:4}\;d_{3:0}\;K_{3:0}
  \qquad\text{where } d = R_d - 16
\]

Three things fall out of that layout, and all three are facts about the machine rather than about
notation.

\textbf{The immediate is split in half.} Bits 11 to 8 carry the top nibble and bits 3 to 0 the
bottom one, with the register number in between. There is no reason for this beyond the encoding
being easier to decode in hardware that way, and it is the detail that makes hand-assembling an
\code{ldi} feel like a puzzle the first time.

\textbf{The register field is four bits, so it counts 0 to 15}, and the machine adds 16. That is
the whole explanation for \secref{c1:sec:regfile}'s restriction: \code{ldi} cannot name \code{r15}
because there is nowhere in the instruction to put the number 15 and mean the sixteenth register.

\textbf{The opcode is only four bits.} \code{1110} means \code{ldi} and nothing else, which is a
lot of encoding space for one instruction, and is why the immediate can afford to be a full byte.

So \code{ldi r16, 0x2A} is \code{1110 0010 0000 1010} = \code{0xE20A}. You will confirm that
number three ways in this chapter: by hand, in the listing \code{avra} writes beside your source,
and by looking at what \code{avr-objdump} says the assembler produced (\secref{c1:sec:disasm}).

\subsection{What this machine does not have}
\label{c1:sec:absent}
Naming the absences early saves you looking for them later.

\textbf{No barrel shifter.} \code{lsl r18} shifts by exactly one bit. There is no instruction that
shifts by a variable number of places, so \code{1 << n} for a runtime \code{n} is a \emph{loop},
and it costs about six cycles per bit. That is the whole reason \code{shift_bits}
(\secref{c1:sec:shiftbits}) exists, and it is worth sitting with: an operation that is free in C is
a loop here, and a driver that calls it with a pin number costs different amounts for different
pins.

\textbf{No division.} There is a hardware multiply (\code{mul}, two cycles), but division is a
library routine, and a slow one. Powers of two are shifts; everything else is expensive.

\textbf{No stack frame, and no stack pointer discipline.} There is a stack pointer, \code{SP}, and
\code{rcall} and \code{ret} use it, and that is the entire mechanism. Local variables, parameter
passing beyond a few registers, and reentrancy are conventions you follow or fail to follow.
Nothing enforces them. Chapter~2 is about the convention; Chapter~4 is about the stack itself.

\textbf{No protection of any kind.} Writing past the end of an array walks into whatever is next.
Pointing \code{Z} at \code{0x0000} and storing writes to \code{r0}. Recursing too deep grows the
stack down into your variables and there is no fault, no message, and no crash, just wrong
answers.

\textbf{And no \code{nop}-free timing.} There is no cache, no branch predictor and no pipeline
stall you have to reason about, which is the good news hiding inside all of the above: an
instruction's cost is a number in a table, the same number every time. That is what makes
\secref{c1:sec:cycles} possible, and it is a property most machines lost decades ago.

% Section 1.2, from lectures/L01/appendix/b_toolchain.md.

\section{The Toolchain}
\label{c1:sec:toolchain}

\subsection{What assembling actually does}
\label{c1:sec:assembling}
An assembler is the least clever program in your toolchain. It reads one line at a time, turns
each instruction into the one or two words that encode it (\secref{c1:sec:encoding}), resolves the
labels, and writes the result out. It does not reorder anything, it does not allocate registers,
and it does not check that what you wrote makes sense. If you write \code{ldi r16, 0x2A}, you get
\code{0xE20A}, and if you write a loop with no exit you get a loop with no exit.

That is the whole appeal. Everything the machine does is something you wrote.

\bookfigure{toolchain}{The toolchain: \code{led.asm} through \code{avra} into a listing, a hex and
a map. The hex and the map together feed the test suite; the hex alone feeds
\code{avrdude}.}{fig:toolchain}

\subsection{\code{avra} as an assembler}
\label{c1:sec:avra}
This book assembles with \code{avra}, which takes AVRASM2, the syntax Microchip Studio assembles,
and never uses the C compiler except in Chapter~6:

\begin{shell}
avra -I /usr/share/avra -I drivers/include -D F_CPU=16000000 \
     -l drivers/build/drivers.lst -m drivers/build/drivers.map \
     -o drivers/build/drivers.hex drivers/build/drivers_unit.asm
\end{shell}

Each of those is load-bearing:

\begin{booktable}{@{}p{10.4em}L@{}}
  \thead{Flag} & \thead{Why} \\
  \midrule
  \mbox{\code{-I /usr/share/avra}} & Where \code{m328Pdef.inc} lives, which is what gives you
    \code{PORTB} and \code{RAMEND}. \\
  \mbox{\code{-I drivers/include}} & Where the course's own \code{.inc} contracts live. \\
  \mbox{\code{-D F_CPU=16000000}} & The clock, as a symbol your own source may use. Nothing in the book
    requires it; Chapter~5's arithmetic is easier to write against a name than a literal. \\
  \code{-l ....lst} & The listing: your source with the encoding beside every line.
    \secref{c1:sec:disasm}. \\
  \code{-m ....map} & The map: every label and its address. \textbf{The tests read this.} \\
  \code{-o ....hex} & Intel hex, which is what the simulator loads. \\
\end{booktable}

You will rarely type that, because \code{make build} does it. Two of the values are overridable
from the environment for the same reason: \code{F_CPU=8000000 make build} assembles against a
different clock, and \code{AVRA_INCLUDE=... make build} is what a non-apt install of \code{avra}
needs, because its device files will not be in \file{/usr/share/avra}.

\textbf{The device is chosen by the file you include, not by a flag.} There is no \code{-mmcu}
here; \code{.include "m328Pdef.inc"} is the whole of that decision, and it is why the line belongs
at the top of every source file you write.

\textbf{\code{avra} has no linker.} It assembles exactly one file, so \file{ci/build.sh} writes a
top-level unit that \code{.include}s every source in \file{drivers/source} and assembles that
instead. You never edit it. When you want to know what was assembled and in what order, read
\file{drivers/build/drivers_unit.asm}; it is one line per source plus two, and generated fresh on
every build.

Two conventions the files in this book follow:
\begin{itemize}
  \item \textbf{\code{;} starts a comment}, to the end of the line.
  \item \textbf{Every label is global already.} \code{avra} exports all of them, so there is no
    \code{.global} to write and no way to forget it. A test that cannot find your subroutine is
    telling you about a spelling or a file that did not assemble, never about a missing
    directive.
\end{itemize}

\subsection{Reading a disassembly}
\label{c1:sec:disasm}
The listing file is the assembler telling you what it produced, line by line, against the source
you gave it. \code{make build} writes one every time:

\begin{shell}
less drivers/build/drivers.lst
\end{shell}

\begin{console}
C:000000 01fc          movw r30, r24            ; Z = the struct
C:000001 e203          ldi  r16, low(PINB + 0x20)
C:000002 8300          std  Z + LED_PIN_REG, r16
\end{console}

The first column is the \textbf{word} address, the second is the encoding, and the rest is your
own line, comment included. That is the source you wrote, the bits it became, and the address
each landed at, all on one line.

To disassemble instead, starting from the machine code and recovering the instructions, which is
the skill this book is actually after, point \code{avr-objdump} at the hex:

\begin{shell}
avr-objdump -D -m avr:5 -b ihex drivers/build/drivers.hex
\end{shell}

\code{-b ihex} is needed because a hex file says nothing about what architecture it is for, and
\code{-m avr:5} supplies the answer. \textbf{The symbol names are gone}, and the output says
\code{.sec1} where an ELF would have said \code{<led_init>}; a hex file carries addresses and
bytes and nothing else, which is exactly why the tests read the map file alongside it.

\begin{console}
Disassembly of section .sec1:

00000000 <.sec1>:
   0:   21 e0           ldi     r18, 0x01       ; 1
   2:   30 e0           ldi     r19, 0x00       ; 0
   4:   38 17           cp      r19, r24
   6:   19 f0           breq    .+6             ; 0xe
   8:   22 0f           add     r18, r18
   a:   33 95           inc     r19
   c:   fb cf           rjmp    .-10            ; 0x4
   e:   82 2f           mov     r24, r18
  10:   08 95           ret
\end{console}

Four columns: the byte address, the encoded words, the instruction, and a comment. Four things in
that listing are worth noticing on your first read.

\textbf{There is one block and it is called \code{.sec1}.} Not \code{shift_bits}, not
\code{shift_bits_loop}: a hex file has no symbol table, so \code{objdump} invents one section name
for the whole run of bytes and labels every branch target with a bare address. The names are in
\file{drivers.map}, which is the other file \code{avra} wrote and the one the tests read. Compare
this against the listing above, where your own label sits on the line that produced it. The
listing knows the names because it is the assembler talking, and the disassembly does not because
it is only ever given bytes.

\textbf{The bytes are little-endian.} \code{ldi r18, 0x01} shows as \code{21 e0}, and the word is
\code{0xE021}. Hold that against \secref{c1:sec:encoding}: opcode \code{1110}, \code{K[7:4]} =
\code{0000}, \code{d} = \code{0010} (so \code{r18}), \code{K[3:0]} = \code{0001}. It checks out,
and reading it in the wrong byte order is a small rite of passage.

\textbf{Addresses here are bytes, not words.} The \code{rjmp} at the bottom goes back to
\code{0x04}, which is word \code{0x02} and so the \emph{third} instruction: words \code{0x00} and
\code{0x01} hold the two \code{ldi}s above it. The vector table is quoted in words and the
disassembly in bytes, and keeping the two apart is a recurring nuisance that this section cannot
make go away.

\textbf{\code{lsl r18} came back as \code{add r18, r18}.} They are the same instruction: shifting
left by one and adding a number to itself do the same thing to the bits and set the same flags,
so the assembler encodes \code{lsl} as \code{add} and the disassembler has no way to know which
you wrote. Several AVR mnemonics are aliases like this, and a disassembly that does not match your
source line for line is usually one of them rather than a mistake.

\subsection{Two images, for two jobs}
\label{c1:sec:images}
\code{make build} produces two, and confusing them wastes an afternoon.

\textbf{\file{drivers/build/drivers.hex}} is the subroutines alone. It has no reset vector and no
vector table; nothing ever runs it from the start. The unit tests set the program counter
straight to a symbol and call one subroutine at a time. This is the image nearly everything in
this book loads, and it is the one \code{make measure} uses by default.

\textbf{\file{drivers/build/app.hex}} is \file{drivers/app/main.asm} assembled together with
those same subroutines: a whole program with a vector table and a main loop. This chapter's
program test and Chapter~3's integration test load this one and let it run, because a program
cannot be called and an interrupt has to arrive. Reach it with \code{make measure IMAGE=app}.

Each comes with a \code{.map} beside it, and the two are read together: the hex is the program and
the map is the names. Delete one and the harness reports the other missing.

\subsection{Running the tests}
\label{c1:sec:tests}
From the repository root:

\begin{shell}
make test
\end{shell}

Every chapter's suite is cumulative, so this runs Chapter~1's now and will still be running it in
Chapter~5. What is in a suite decides for itself whether to run: the tests that check the pinned
device constants always run, and the ones that call your assembly are switched on by their
\code{.asm} file existing. Nothing has to be registered anywhere; the README beside the suite, in
\file{lectures/L01/exercises/test}, has the details.

Two failures worth recognising before you meet them:
\begin{itemize}
  \item \textbf{\code{Utils.SubroutinesAreDefined} fails and the file is plainly there.} The label
    is spelled differently from the name the test asks for, or the file did not assemble and the
    build said so further up. \code{avra} needs no \code{.global}, so it is never that.
  \item \textbf{A cycle count comes back as 100001.} That is the simulator's budget running out,
    which is what an infinite loop looks like from outside.
\end{itemize}

\subsection{Measuring a subroutine}
\label{c1:sec:measure}
A test tells you pass or fail. A cross-check needs the number itself, so the companion repository
ships a tool:

\begin{shell}
make measure SYMBOL=shift_bits ARG=5
\end{shell}

\begin{console}
shift_bits(r25:r24 = 0x0005, r22 = 0)
  cycles     40
  time       2.500 us at 16.0 MHz
  returned   r24 = 32 (0x20)

deepest stack: SP reached 0x08FD, 2 bytes below RAMEND
\end{console}

It assembles whatever is in \file{drivers/source}, loads it, sets \code{r25:r24} and \code{r22}
to the arguments you gave, calls the symbol, and reports what came back and what it cost.
\code{IMAGE=app} measures against the whole program instead of the library, which is what an
interrupt handler needs (\secref{c1:sec:images}).

\textbf{The figure excludes the \code{rcall} that would have got you there}, because there was no
\code{rcall}: the tool sets the program counter directly, the way the tests do. That is not a
defect and it is not hidden from you; it is the first term of the discrepancy that
Exercise~\ref{c1:ex:crosscheck}, the cross-check, asks you to account for.

\subsection{Opening this book's code in Microchip Studio}
\label{c1:sec:studio}
Everything in \file{drivers/source} is AVRASM2, which is what Microchip Studio's assembler takes.
A file from this book opens there and single-steps without being translated first, and that is
the reason for the choice: the Processor Status window showing \code{R16} change, and the I/O
view letting you tick a bit of \code{PINB} to fake a button press, apply to \textbf{your} code
rather than to a port of it.

Nothing in this book requires Studio, and no exercise depends on it. It is Windows-only, and
\code{make test} measures everything Studio would show you and rather more besides. But if you
have it, this is how it fits:
\begin{enumerate}
  \item \textbf{File $\rightarrow$ New $\rightarrow$ Project $\rightarrow$ AVR Assembler
    Project}, targeting the ATmega328P.
  \item Replace the generated \code{.asm} with one of yours, or add yours to the project.
  \item \textbf{Debug $\rightarrow$ Start Debugging and Break}, and pick \textbf{Simulator} as the
    tool when asked.
  \item \textbf{Debug $\rightarrow$ Windows $\rightarrow$ Processor Status} for the register file,
    and \textbf{I/O} for the peripherals.
\end{enumerate}

Two differences from the build here, neither of which is a syntax question:
\begin{itemize}
  \item \textbf{Studio assembles one file and this repository assembles all of them.}
    \file{ci/build.sh} generates a unit that includes every source (\secref{c1:sec:avra}); a Studio
    project has whatever you put in it. A subroutine that assembles here and is missing there is
    usually a file you did not add.
  \item \textbf{\code{m328Pdef.inc} comes from a different place.} Studio ships its own copy and
    finds it automatically; \code{avra} uses the one in \file{/usr/share/avra}. They agree about
    the ATmega328P, which is the only device this book targets.
\end{itemize}

\subsubsection{The vector table, in the units the datasheet uses}
\code{.org} counts \textbf{words} in AVRASM2, and so does the datasheet's vector table.
\code{.org 0x0006} is PCINT0 in both, and no conversion stands between them. Better still,
\code{m328Pdef.inc} names them:

\begin{asm}
.org PCI0addr
    rjmp isr_pcint0
\end{asm}

\code{PCI0addr}, \code{WDTaddr}, \code{OC1Aaddr} and the rest are defined in the device file, so
the magic number need not appear in your program at all. Chapter~3 is where this matters.

\begin{aside}
\textbf{Reading GNU \code{as} code elsewhere?} Most AVR assembly published online is written for
\code{avr-gcc}, where \code{.equ} takes a comma, register names come from \code{<avr/io.h>},
\code{PORTB} means the data space address \code{0x25} rather than the I/O address \code{0x05},
and \textbf{\code{.org} counts bytes}, so the same vector table is written \code{.org 0x000C}.
\secref{c6:sec:asmfromc} shows that dialect properly, because Chapter~6 is the one chapter where this book links
against a C compiler.
\end{aside}

% Section 1.3, from lectures/L01/appendix/c_counting_cycles.md.

\section{Counting Cycles}
\label{c1:sec:cycles}

\subsection{Why bother}
\label{c1:sec:whybother}
On most machines you cannot say what a routine costs. Caches, branch prediction and out-of-order
execution mean the same code takes different times on different runs, and the honest answer to
``how long does this take'' is ``measure it, repeatedly, and take a distribution''.

The ATmega328P has none of those. Every instruction takes a fixed number of cycles, the number is
in a table, and it is the same number every time. So ``how long does this take'' has an exact
answer you can work out on paper before the code exists, and that changes what is worth doing:
\begin{itemize}
  \item You can decide between two implementations without writing either.
  \item You can tell whether a delay loop will meet its deadline without a scope.
  \item When a measurement disagrees with your count, one of you is wrong about the code, and
    finding out which is a productive afternoon rather than a guess about the hardware.
\end{itemize}

This section is that table, and the arithmetic on top of it.

\subsection{What each instruction costs}
\label{c1:sec:costs}
Every figure here comes from the AVR Instruction Set Manual for this core, and every one was
confirmed by running that instruction in the simulator and subtracting the \code{ret} around it.
Where the two had disagreed, the disagreement would be written down here. They did not.

\begin{booktable}{@{}p{3.6em}L@{}}
  \thead{Cycles} & \thead{Instructions} \\
  \midrule
  1 & \code{ldi mov movw add adc sub subi and andi or ori eor com neg cp cpc cpi lsl lsr rol ror
      asr inc dec nop in out sei cli} \\
  2 & \code{rjmp lds sts ld ldd st std push pop adiw sbiw sbi cbi} \\
  3 & \code{rcall lpm} \\
  4 & \code{ret reti} \\
  1 or 2 & \code{breq brne brlo brsh brcc brcs} and the other conditional branches \\
\end{booktable}

Three rows in that table are worth more attention than the rest.

\textbf{\code{ret} costs four cycles and \code{rcall} costs three.} Calling a subroutine that does
nothing at all costs seven cycles, which is more than most loops' bodies. That is not an argument
against subroutines; it is the reason a driver's cost is never just the cost of its body, and the
reason the cross-check in Exercise~\ref{c1:ex:crosscheck} comes out the way it does.

\textbf{\code{sbi} and \code{cbi} cost two cycles, not one.} They look like single-bit versions of
\code{out}, and they are not: they read the register, modify one bit, and write it back. Worth
remembering when you are counting a routine that sets several bits one at a time and wondering
where the cycles went.

\textbf{A memory access costs twice what a register operation does.} \code{lds} is two cycles and
\code{mov} is one, which is the whole argument for keeping working values in registers. On a
machine with 32 of them that is usually possible, and arranging for it is most of what makes
assembly fast.

\subsection{A branch costs more when it branches}
\label{c1:sec:branchcost}
A conditional branch takes \textbf{one cycle if it falls through and two if it is taken}. The
reason is that a taken branch changes the program counter, and the instruction already being
fetched has to be thrown away.

This is the single most common source of an off-by-a-few cycle count, because it means a loop's
cost is not the number of iterations times the body.

Take a loop with its test at the top, which is the shape the subroutine in \secref{c1:sec:first}
needs: compare a running count against the argument, branch out of the loop when they are equal,
do the work, add one to the count, and jump back to the top.

\begin{booktable}{@{}Lp{10.6em}l@{}}
  \thead{Step} & \thead{Cycles} & \thead{Times it runs} \\
  \midrule
  compare the count against the argument & 1 & $n + 1$ \\
  branch out when they are equal & 1 falling through, 2 taken & $n + 1$ \\
  the work, one instruction here & 1 & $n$ \\
  add one to the count & 1 & $n$ \\
  jump back to the top & 2 & $n$ \\
\end{booktable}

A loop running $n$ times executes the branch \textbf{$n + 1$ times}: $n$ of them fall
through into the body, and one is taken to leave. So the branch contributes
$n \times 1 + 1 \times 2 = n + 2$ cycles, not $n + 1$ and not $2n$.

Everything else in the loop runs exactly $n$ times, except the comparison, which also runs on the
final pass that leaves the loop and so runs $n + 1$ times as well. That extra comparison is the one
people forget, and it is why the answer comes out ten rather than nine when $n$ is zero.

\subsection{Counting a routine on paper}
\label{c1:sec:onpaper}
The table in \secref{c1:sec:branchcost} is the method, and it is worth stating as a method because
you will use it in every chapter of this book.

\textbf{One row per instruction, not per line executed.} Write the instruction, its cost from
\secref{c1:sec:costs}, and how many times it runs. Multiply and sum. A branch gets two rows, because
it costs differently on the passes that fall through and on the one that is taken.

For \code{shift_bits} with an argument of $n$:

\begin{narrowtable}{@{}lccc@{}}
  \thead{Instruction} & \thead{Cycles} & \thead{Times} & \thead{Total} \\
  \midrule
  \code{ldi} $\times$ 2 & 1 & 2 & 2 \\
  \code{cp} & 1 & $n + 1$ & $n + 1$ \\
  \code{breq}, falling through & 1 & $n$ & $n$ \\
  \code{breq}, taken & 2 & 1 & 2 \\
  \code{lsl} & 1 & $n$ & $n$ \\
  \code{inc} & 1 & $n$ & $n$ \\
  \code{rjmp} & 2 & $n$ & $2n$ \\
  \code{mov} & 1 & 1 & 1 \\
  \code{ret} & 4 & 1 & 4 \\
\end{narrowtable}

\noindent which adds to $6n + 10$.

\textbf{Order does not matter and that is not a simplification.} Cost is a sum, so a table that
tracked which instruction ran when would be carrying weight for nothing. What the method does
require is that you have already worked out \textbf{how many times each line runs}, and that is
the part that is actually hard. \Secref{c1:sec:branchcost} is about exactly that, and the $n + 1$ on
the comparison is where most wrong answers come from.

\subsection{What a hand count is blind to}
\label{c1:sec:blind}
It is worth being explicit, because a method that quietly stops being right is worse than no
method.

\textbf{It does not know about the \code{rcall} that reached the routine.} Your table is the body.
The three cycles of getting there, and the instructions that set up the arguments, belong to the
caller, and counting them is the caller's table's job.

\textbf{It does not know about interrupts.} An interrupt arriving mid-routine adds the handler's
cost plus six to nine cycles of entry and four of \code{reti}, at a moment nothing in this model can
predict. From Chapter~3 onwards that is a real effect and a hand count will simply be wrong about
it, by an amount that depends on what else is running. Where that matters, the answer is to
measure.

\textbf{It does not know whether your table matches your code.} Nothing checks that the rows you
wrote describe the routine you assembled. Two things guard against that: the simulator, which
costs the real thing, and the discipline of writing the table from the source rather than from
memory. When they disagree, the table is usually the one that is wrong, because it was written by
a human who was thinking about something else.

\textbf{It says nothing about correctness.} A routine can cost exactly what you predicted and
return entirely the wrong answer. Cycle counting tells you whether something fits in the time
available, and nothing whatever about whether it should be running at all.

% Section 1.4, from lectures/L01/appendix/d_first_subroutine.md.

\section{Your First Subroutine}
\label{c1:sec:first}

\subsection{The task}
\label{c1:sec:task}
Write two subroutines in \file{drivers/source/utils.asm}, and one program in
\file{drivers/app/main.asm}.

The subroutines compute \code{1 << n} and \code{~(1 << n)} for a value of \code{n} that is not
known until the program runs. In C that is one operator each and you would not give it a second
thought. Here it is a loop, because the AVR has no barrel shifter (\secref{c1:sec:absent}):
\code{lsl} shifts by exactly one place, and shifting by \code{n} places means going round \code{n}
times.

That is not a curiosity. Every driver in this book is handed a pin number and has to turn it into
a bit mask, and these two subroutines are how. \code{led_on} in Chapter~2 calls
\code{shift_bits}; \code{btn_init} in Chapter~3 calls both. Getting them right now means the rest
of the book is about peripherals rather than about bit twiddling.

Creating the file is all it takes to switch its tests on; the README beside the suite, in
\file{lectures/L01/exercises/test}, says how.

\subsection{Where the arguments live}
\label{c1:sec:args}
This book follows the AVR C calling convention from the first line of assembly it writes, and
keeps following it, so that Chapter~6's ``call your assembly from C'' is a demonstration rather
than a rewrite. The whole of what you need for now:

\begin{booktable}{@{}p{9.4em}L@{}}
  \thead{Register} & \thead{Role} \\
  \midrule
  \code{r24} & First argument, if it is a byte. Also where a byte-sized return value goes. \\
  \code{r25:r24} & First argument, if it is a pointer or another 16-bit value. \\
  \code{r22} & Second argument, if it is a byte. \\
  \code{r18} to \code{r27}, \code{r30}, \code{r31} & \textbf{Call-clobbered.} A subroutine may
    destroy these freely. \\
  \code{r2} to \code{r17}, \code{r28}, \code{r29} & \textbf{Call-saved.} A subroutine that uses
    one must put it back. \\
\end{booktable}

So both subroutines here take their argument in \code{r24} and leave their result in \code{r24},
and both may use \code{r18} and \code{r19} as scratch without saving anything.

\subsection{\code{shift_bits}}
\label{c1:sec:shiftbits}
Returns \code{1} shifted left by the number of places given.

\begin{narrowtable}{@{}ccl@{}}
  \thead{Given \code{r24}} & \thead{Returns \code{r24}} & \thead{Because} \\
  \midrule
  0 & \code{0x01} & Shifted no times at all. \\
  1 & \code{0x02} & \\
  5 & \code{0x20} & Which is bit 5, the Arduino Uno's pin 13. \\
  7 & \code{0x80} & \\
  8 & \code{0x00} & The bit has been shifted off the end of an eight-bit register. \\
\end{narrowtable}

Write it like this:
\begin{enumerate}
  \item Put \code{0x01} in a scratch register; this is the value being shifted. Put \code{0} in a
    second scratch register; this is the count of shifts done so far.
  \item At the top of the loop, compare the count against \code{r24}. If they are equal, leave the
    loop.
  \item Otherwise shift the value left once, add one to the count, and jump back to the top.
  \item On leaving, move the value into \code{r24} and return.
\end{enumerate}

Two things about that shape are worth being deliberate about, because they are what the cost
depends on (\secref{c1:sec:branchcost}):

\textbf{The comparison is at the top, not the bottom.} So an argument of zero does no shifts at
all, which is what makes \code{shift_bits(0)} return \code{0x01} rather than \code{0x02}. A loop
with the test at the bottom always runs its body once, and would be wrong for exactly one input,
which is the input your driver will use for pin 8.

\textbf{The exit branch runs once more than the body does.} That last comparison, the one that
finds the count has reached \code{r24} and leaves, is a real instruction with a real cost. It is
the reason the answer is 10 cycles at $n = 0$ and not 7.

The last row of the table is the interesting one. Nothing in the routine says ``stop at eight'';
the zero falls out of the register being eight bits wide, and the bit simply leaves. A reader
thinking in C, where \code{1 << 8} is 256, will predict \code{0x100} and there is nowhere to put
it.

\subsection{\code{shift_bits_inverted}}
\label{c1:sec:inverted}
Returns the complement of the same value: \code{~(1 << n)}.

\begin{narrowtable}{@{}cc@{}}
  \thead{Given \code{r24}} & \thead{Returns \code{r24}} \\
  \midrule
  0 & \code{0xFE} \\
  1 & \code{0xFD} \\
  5 & \code{0xDF} \\
  7 & \code{0x7F} \\
  8 & \code{0xFF} \\
\end{narrowtable}

The obvious implementation is to call \code{shift_bits} and complement the result with
\code{com}. Do not; write it as its own loop. Start the scratch value at \code{0xFE} instead of
\code{0x01}, and each time round shift left \emph{and then increment}, which shifts a zero in at
the bottom and turns it back into a one.

The reason for the detour is that this routine exists to produce a mask for clearing a bit, and
Chapter~2's driver calls it directly. Writing it as its own loop also makes the point that
\secref{c1:sec:branchcost} is really making: two routines differing by a single \code{inc} differ in
cost by that \code{inc} on every trip, and you can say by how much before running either.

\subsection{Why \code{r18} and \code{r19}}
\label{c1:sec:whyr18}
Use \code{r18} and \code{r19} as the scratch registers, not \code{r16} and \code{r17}.

Both pairs are in the upper half, so both can take an \code{ldi} (\secref{c1:sec:regfile}). The
difference is the calling convention in \secref{c1:sec:args}: \code{r18} and \code{r19} are
call-clobbered, and \code{r16} and \code{r17} are call-saved. A subroutine that writes to
\code{r16} without restoring it has broken its contract with whoever called it, and the caller has
no way to know.

Right now nothing calls these subroutines but a test, and a test does not care. By Chapter~6 the
caller may be a C function that had a live value in \code{r16}, and then the bug is a variable
that changes on its own across a function call, which is about as unpleasant a thing to debug as
this book contains. Choosing the right registers now costs nothing and means that never happens.

If you are comparing against AVR assembly you have found elsewhere, note that a good deal of it
uses \code{r16} and \code{r17} here. It works, right up until it does not.

\subsection{The program}
\label{c1:sec:program}
\file{drivers/app/main.asm} is the first thing in this book that is a program rather than a
subroutine. For this chapter it needs four things:

\textbf{A reset vector.} The first word of program memory is where the machine starts, so the
first instruction in your file has to be a jump to your code. Use \code{.org 0} followed by an
\code{rjmp}.

\textbf{A stack pointer.} \code{rcall} pushes a return address and \code{ret} pops one, and both
use \code{SP}. Set it to \code{RAMEND} by writing \code{high(RAMEND)} to \code{SPH} and then
\code{low(RAMEND)} to \code{SPL}, with \code{out} rather than \code{sts}
(\secref{c1:sec:iowindow}). \code{low()} and \code{high()} are AVRASM2's; GNU \code{as} spells the
same two functions \code{lo8()} and \code{hi8()}, which matters once in Chapter~6 and nowhere
else.

The hardware already does this on the ATmega328P, and it does it after \textbf{every} reset,
watchdog resets included. You will find a good deal of AVR code carrying a comment saying
otherwise; on this device that comment is wrong, and repeating it is how the belief survives.

Write the four instructions anyway, for the two reasons that are actually true. Older AVRs reset
\code{SP} to \code{0x0000} rather than to \code{RAMEND}, so code that omits this is not portable
to them, and the failure there is a stack that grows down out of memory on the first
\code{rcall}. And a reset is not the only way to arrive at your code: a bootloader that jumps into
it has performed no reset at all, and \code{SP} holds whatever the bootloader left in it.

Neither of those applies to anything in this book. Knowing \emph{which} reason applies, rather
than carrying a habit you cannot justify, is the point.

\textbf{Something to do, and somewhere to stop.} Call \code{shift_bits} with an argument, then sit
in a loop that jumps to itself. There is no operating system to return to; a program that runs off
the end of its own code executes whatever is in flash after it, which is \code{0xFF} repeated, and
the results are not interesting. An infinite loop is how an embedded program ends.

\textbf{And something you can see.} Light the LED on PB5.

\begin{asm}
    sbi  DDRB, 5                ; PB5 is an output
    sbi  PORTB, 5               ; and it is high
\end{asm}

That is the whole of it, and every part of it is already in this chapter. \code{DDRB} is
\code{0x04} and \code{PORTB} is \code{0x05}, both I/O addresses, so \code{sbi} reaches them and
\code{sts} would not (\secref{c1:sec:iowindow}); \code{sbi} costs two cycles and not one, which
\secref{c1:sec:costs} says and which is worth noticing now rather than in Chapter~5 when you are
counting a timer's handler.

The one thing that is new is what the two registers mean, and it is one sentence:
\textbf{\code{DDRx} decides whether a pin is an output, and \code{PORTx} decides what an output
drives.} Both are per-bit, both default to zero, and zero means ``input, not driven''. So an LED
needs both writes and in that order: setting \code{PORTB} first would briefly enable a pull-up on
an input pin instead, which is a different thing that happens to look similar.

PB5 is Arduino pin 13, which is the one with an LED already fitted to the board, and the reason
this book lights that one rather than any other. On a real Uno there is a resistor beside it; in
the simulator there is nothing to protect, and \secref{c2:sec:resistor} is where that stops being a detail.

\textbf{Two instructions, one pin, hard-coded.} That is deliberate, and it is exactly as far as
this chapter goes. It only works for PB5, it cannot be told which LED to light, and nothing else
in the program can reuse it. Turning it into something a program can call with a pin number is
Chapter~2, and so is the reason that turns out to be more work than it sounds.

\subsection{How to check it}
\label{c1:sec:check}
Three things, in this order.

\textbf{Does it assemble?}

\begin{shell}
make build
\end{shell}

\textbf{Does it do the right thing, and cost the right amount?}

\begin{shell}
make test
\end{shell}

Nine tests were passing before you wrote anything, eighteen once \file{utils.asm} exists, and
twenty once \file{main.asm} does. If \code{Utils.SubroutinesAreDefined} fails while the file is
plainly there, the label is spelled differently from the name the test asks for, or the file did
not assemble (\secref{c1:sec:avra}).

The tests for \file{main.asm} work differently from the rest, and it is worth knowing why before
one of them fails. A subroutine can be called: a test sets the program counter to
\code{shift_bits}, puts a number in \code{r24} and reads \code{r24} back. A program cannot. It has
no arguments, it never returns, and the only way to find out what it does is to \textbf{let it
run} and then look at the machine:

\begin{cppcode}
mcu.run(2000);                      // 2000 cycles of your program
mcu.data(0x24) & (1 << 5)           // DDRB, in the data space
\end{cppcode}

Note the \code{0x24}. The test reads the \textbf{data space} address where your program wrote the
\textbf{I/O} address \code{0x04}, and both are correct, because they are two names for one
register. That is \secref{c1:sec:iowindow} again, seen from the third side, and it is the last time
this book will point it out.

\textbf{What does it actually cost?}

\begin{shell}
make measure SYMBOL=shift_bits ARG=5
\end{shell}

Do not run that until you have predicted the answer. Predicting it by hand, and then reconciling
the prediction against the measurement, is the cross-check in Exercise~\ref{c1:ex:crosscheck}, and
it is the exercise this whole chapter exists to set up.

% Section 1.5, from lectures/L01/README.md: what you should be able to do, questions to test
% yourself, the next lecture, and the references.

\section{Review}
\label{c1:sec:review}

\needspace{6\baselineskip}
\subsection*{What you should be able to do now}
\begin{itemize}
  \item Say what the 32 registers are, which sixteen an immediate instruction can name, and why
    the other sixteen cannot be named that way.
  \item Read a line of \code{SREG}'s contents and say which branch would be taken next.
  \item Given an address, say which of the three address spaces it belongs to and which
    instruction reaches it; and say why \code{in}, \code{lds} and \code{lpm} are three different
    instructions rather than one.
  \item Encode an \code{ldi} by hand and confirm it against a disassembly.
  \item Assemble a \code{.asm} file, disassemble the result, and run it in the simulator.
  \item Count the cycles of a straight-line routine and of a loop, from the instruction set
    summary, and say why a measured figure at a call site comes out higher than your count of the
    body.
  \item Write a subroutine that takes an argument in \code{r24}, returns a value in \code{r24},
    and returns.
  \item Write a program that runs on its own: a reset vector, a stack pointer, an LED lit through
    \code{DDRB} and \code{PORTB}, and a loop to stop in.
\end{itemize}

\needspace{6\baselineskip}
\subsection*{Questions to test yourself}
\begin{enumerate}
  \item Why can \code{ldi} name \code{r16} but not \code{r15}, and what does that cost you when you
    run out of registers in the upper half?
  \item Data space address \code{0x0025} and I/O address \code{0x05} are the same register. Which
    instructions take which of the two numbers, and what happens if you give one the other's?
  \item The register file is addressable as data at \code{0x0000} to \code{0x001F}. What could you
    do with that, and why does essentially no program do it?
  \item A conditional branch costs one cycle when it falls through and two when it branches. A
    loop running $n$ times executes its exit branch $n + 1$ times. How many of those are taken?
  \item Your hand count of a subroutine says 40 cycles and the simulator agrees, yet the
    \emph{call} costs 44 and the loop around it costs 46 per iteration. Where are the other four,
    and the two after that, without looking anything up?
  \item Why does \code{shift_bits(7)} cost more than five times what \code{shift_bits(0)} costs,
    and what does that imply about a driver that calls it with a pin number?
  \item Your program writes \code{DDRB} at I/O address \code{0x04}, and the test that checks it
    reads data space address \code{0x24}. Both are right. Which instructions take which of the two
    numbers?
  \item You light an LED with \code{sbi PORTB, 5} and it works. Why can a driver that has to light
    \emph{any} LED not be written that way at all?
\end{enumerate}

\needspace{6\baselineskip}
\subsection*{Looking ahead}
Chapter~2 takes the same LED and lights it with something that can be told which LED to light:
\begin{itemize}
  \item \code{DDRx}, \code{PORTx} and \code{PINx}: three registers per port, and what each of the
    three actually does, which is more than the two \code{sbi} instructions at the end of this
    chapter needed to know.
  \item Why writing a one to \code{PINx} toggles an output, which is not what the name suggests at
    all.
  \item The calling contract: which registers a subroutine may destroy, and which it must give
    back.
  \item The LED driver: the same LED, lit by something that can be told which LED to light.
\end{itemize}

\needspace{6\baselineskip}
\subsection*{Sources}
Every number this chapter quotes comes from one of two primary documents, and the text says
which:
\begin{itemize}
  \item \emph{ATmega328P Datasheet}, Microchip Technology.\\
    \url{https://ww1.microchip.com/downloads/en/DeviceDoc/Atmel-7810-Automotive-Microcontrollers-ATmega328P_Datasheet.pdf}
  \item \emph{AVR Instruction Set Manual}, Microchip Technology.\\
    \url{https://ww1.microchip.com/downloads/en/devicedoc/atmel-0856-avr-instruction-set-manual.pdf}
\end{itemize}

% Section 1.6, from lectures/L01/appendix/e_exercises.md.

\section{Exercises}
\label{c1:sec:exercises}

\begin{aside}
\textbf{How to check your work.} Exercises \ref{c1:ex:firstsub} and \ref{c1:ex:program} are checked by
this chapter's test suite: write the file at the path the specification gives, then run
\code{make test} from the repository root. Nothing has to be registered; creating the file is what
switches its tests on.

The rest will have worked solutions in the companion repository, published later as
\file{lectures/L01/appendix/f_solutions.md}, and the hand calculations have short answers in
Appendix~\ref{app:answers}. Work each one before you read it. Where an exercise
has a plausible wrong answer, the solution gives that too, and says why
it is wrong, because being told the right answer teaches you less than being shown the one you
nearly gave.

\textbf{No hardware is needed for any of this.} Everything is written, assembled, simulated and
measured on your laptop.
\end{aside}

Each exercise is labelled with its kind. There is exactly one \textbf{Cross-check} per chapter,
and in this one it is Exercise~\ref{c1:ex:crosscheck}.

% ------------------------------------------------------------------------------------------
\exercise{The upper half}{Recall}
\label{c1:ex:upper}
\begin{parts}
  \item State which registers \code{ldi} can name, and which it cannot.
  \item Explain \emph{why}, from the instruction's encoding rather than from the datasheet's
    table. Your answer should mention a specific number of bits.
  \item You are writing a routine that needs six working registers, and you have used \code{r16}
    to \code{r21}. A seventh value is a constant you need to load. Somebody suggests \code{r10},
    which is free. What does that cost you, in instructions and in cycles, every time you set it
    up?
  \item Name two instructions other than \code{ldi} that have the same restriction, and one
    instruction that does not.
\end{parts}

% ------------------------------------------------------------------------------------------
\exercise{Assembling four instructions by hand}{Hand calculation}
\label{c1:ex:assemble}
Using the encoding in \secref{c1:sec:encoding}, and nothing else, work out the 16-bit word each of
these assembles to. Show the four nibbles for at least the first one.
\begin{parts}
  \item \code{ldi r16, 0x2A}
  \item \code{ldi r24, 0x2A}
  \item \code{ldi r31, 0xFF}
  \item \code{ldi r20, 0x03}
  \item One of these is not a legal instruction: \code{ldi r15, 0x2A}. Say what the assembler does
    with it, and what would go wrong if it silently produced a word anyway.
\end{parts}

\noindent\textbf{Check yourself:} the encoding is \code{1110 KKKK dddd KKKK} with
\code{d = Rd - 16}, and the assembler will settle it. Put the instructions in a \code{.asm} file,
assemble it, and run \code{avr-objdump -D -m avr:5 -b ihex} over the hex; the bytes come out
reversed, for the reason in \secref{c1:sec:disasm}.

% ------------------------------------------------------------------------------------------
\exercise{Which space, which instruction}{Recall}
\label{c1:ex:spaces}
For each of the following, say which address space it names, what is stored there, and which
instruction or instructions can reach it.
\begin{parts}
  \item Data space \code{0x0005}
  \item Data space \code{0x0025}
  \item I/O address \code{0x05}
  \item Data space \code{0x0100}
  \item Program word address \code{0x0006}
  \item EEPROM address \code{0x0000}
  \item Two of the above name the same physical register. Which two, and what happens if you use
    one of the two numbers with an instruction expecting the other?
\end{parts}

% ------------------------------------------------------------------------------------------
\exercise{Choosing the branch}{Design}
\label{c1:ex:branch}
\code{r19} holds a loop counter and \code{r24} holds a limit. Both are unsigned.
\begin{parts}
  \item You want to leave the loop when the counter has reached the limit. Write the two
    instructions, and say which flag the second one reads.
  \item You want to stay in the loop while the counter is strictly below the limit. Write the two
    instructions.
  \item Somebody inserts an \code{inc r19} between the comparison and the branch. Explain what
    breaks, and why the symptom is a branch that goes the wrong way rather than an assembler
    error.
  \item Somebody inserts a \code{mov r20, r19} in the same place instead. Explain why this one is
    harmless, and say where you would look it up rather than guessing.
\end{parts}

% ------------------------------------------------------------------------------------------
\exercise{Your first subroutines}{Code}
\label{c1:ex:firstsub}
Write \code{shift_bits} and \code{shift_bits_inverted} in \file{drivers/source/utils.asm}, to the
specification in \secref{c1:sec:first}.
\begin{parts}
  \item Write both. Use \code{r18} and \code{r19} as scratch, for the reason in
    \secref{c1:sec:whyr18}.
  \item Assemble with \code{make build}, then disassemble your own routine with
    \code{avr-objdump -D -m avr:5 -b ihex drivers/build/drivers.hex} and read it against what you
    wrote. Both flags are required, for the reason in \secref{c1:sec:disasm}: a hex file says nothing
    about its own architecture. At least one line will have come back as a different instruction
    than you typed. Say which, and why that is not a mistake.
  \item Run \code{make test}.
  \item Before running anything else: predict what \code{shift_bits} returns for an argument of 8,
    and say what a reader thinking in C would predict instead.
\end{parts}

% ------------------------------------------------------------------------------------------
\exercise{A program that runs on its own}{Code}
\label{c1:ex:program}
Write \file{drivers/app/main.asm} to the specification in \secref{c1:sec:program}.
\begin{parts}
  \item Write the reset vector, the stack pointer initialisation, a call to \code{shift_bits}, the
    two instructions that light PB5, and the final loop.
  \item \code{make build} should now report that it built \file{app.hex} as well as
    \file{drivers.hex}, and \code{make test} should report two more passing tests than it did
    before.
  \item The hardware sets \code{SP} to \code{RAMEND} at reset by itself, so on a cold start your
    four instructions change nothing. Give the case where they do change something, and say what
    would go wrong without them.
  \item Remove the final loop and describe, without running it, what the machine does when it
    reaches the end of your code. What is in flash after your last instruction?
  \item Swap the two \code{sbi} instructions, so that \code{PORTB} is written before \code{DDRB}.
    The LED still ends up lit. Say what the pin is doing during the one instruction in between,
    and why that is not the same thing as being an output that happens to be low.
  \item You wrote \code{sbi DDRB, 5}, and the test that checks it reads data space address
    \code{0x24}. Both numbers are right. Explain, and give the one instruction that would write the
    same register using \code{0x24} directly.
\end{parts}

% ------------------------------------------------------------------------------------------
\exercise{Counting by hand}{Hand calculation}
\label{c1:ex:counting}
Using only the table in \secref{c1:sec:costs} and the rule in \secref{c1:sec:branchcost}, and without
running anything:
\begin{parts}
  \item Count the cycles \code{shift_bits} takes for an argument of 0. Show your working as a line
    per instruction, with how many times each one runs.
  \item Do the same for an argument of 7.
  \item From those two, write down the general formula in terms of $n$.
  \item How many times does the \code{cp} execute when $n$ is 3? How many times does the
    \code{breq}? How many of those \code{breq} executions are taken? These three numbers are not
    all the same, and getting them confused is the usual reason a hand count comes out wrong.
  \item Do the same for \code{shift_bits_inverted}, and state the difference between the two
    formulas in one sentence about a single instruction.
\end{parts}

% ------------------------------------------------------------------------------------------
\exercise[fillaccent]{What a subroutine costs}{Cross-check}
\label{c1:ex:crosscheck}
\emph{Compute it by hand, measure it, reconcile.}

This is the exercise the chapter exists for. Do the parts in order, and write each answer down
before starting the next; the point is lost if you look at the measurement first.
\begin{parts}
  \item \textbf{By hand.} From the exercise above, you have a formula. Evaluate it for $n = 0$ and
    for $n = 7$, and convert both to microseconds at 16\,MHz. One cycle is 62.5 nanoseconds, and
    doing that conversion in your head once is worth more than a function that does it for you.
  \item \textbf{The table, in full.} Write out the count for $n = 7$ in the shape
    \secref{c1:sec:onpaper} gives: one row per instruction, its cost, how many times it runs, and the
    product. You will need it in (d), and a formula you cannot decompose is a formula you cannot
    find the mistake in.
  \item \textbf{By measurement.}
\begin{shell}
make measure SYMBOL=shift_bits ARG=0
make measure SYMBOL=shift_bits ARG=7
\end{shell}
  \item \textbf{Reconcile the two.} They should agree exactly. If they do not, the table is more
    likely wrong than the simulator, and the row to check first is the \code{cp}, which runs
    $n + 1$ times rather than $n$. Find the mistake before continuing rather than adjusting the
    formula until it matches.
  \item \textbf{Now the call site.} Your \file{main.asm} does not just execute \code{shift_bits};
    it loads an argument and calls it. Count, by hand, the cycles of those two extra instructions,
    and write down what a loop that called \code{shift_bits(7)} over and over would cost \emph{per
    iteration}.
  \item \textbf{The discrepancy.} Your answer to (e) is larger than your answer to (a). State the
    difference in cycles, and state it again as a percentage of the (a) figure, once for $n = 0$
    and once for $n = 7$. \textbf{The two percentages are very different.} Explain why, in one
    sentence, without using the word ``overhead''.
  \item \textbf{What it means for a driver.} Chapter~2's \code{led_on} calls \code{shift_bits}
    with a pin number. Using your formula, how many times longer does \code{led_on} take for
    Arduino pin 13 than for Arduino pin 8? (Pin 13 is bit 5 of port B; pin 8 is bit 0 of port B.)
    Say whether that difference would matter, and give one situation in which it would.
  \item \code{make measure} reports a figure that does not include the \code{rcall}. Is that a
    defect in the tool? Give the case for and against, in a sentence each, and say what you would
    have it do instead if you disagree with the choice.
\end{parts}

